% =====================================================================
%  Chapter 3 — The Detection Mandate Against the Characterisation Mandate
% =====================================================================

\chapter{The Detection Mandate Against the Characterisation Mandate}
\label{ch:mandates}

Two mandates sit inside the planetary-defence establishment, and they are
funded, staffed and organised separately. The first is to \emph{find}
near-Earth objects. The second is to \emph{characterise} them. Almost all of
the hardware, almost all of the observing time and almost all of the
public attention are on the first side of that line. The second side has no
automated trigger at all.

That is not a criticism of the discovery effort, which has been extraordinarily
successful. It is an observation about where the residual risk sits once the
catalogue is built.

\section{Five clocks, not one}

The phrase "before it is too late" implies a single clock. There are five, and
they run at very different rates.

\begin{table}[!htbp]
\centering
\small
\caption{The warning-time chain, and the part of it a lunar surface station
would touch.}
\label{tab:clocks}
\begin{tabularx}{\textwidth}{@{}P{0.72in}P{1.12in}P{0.87in}Y@{}}
\toprule
\rowcolor{panel}
\textbf{Stage} & \textbf{Instrument} & \textbf{Clock} & \textbf{Lunar help} \\
\midrule
1. Discovery & wide-field optical survey & minutes of alert & partly: near-Sun
access \\
\rowcolor{panel}
2. Orbit and impact probability & astrometry, B-plane & weeks to years & yes: a
second geometry \\
3. Characterisation & spectra, radar shape, rotation, mass & the bottleneck &
yes, and most \\
\rowcolor{panel}
4. Mitigation decision & modelling & years & no \\
5. Deflection & spacecraft & years & no \\
\bottomrule
\end{tabularx}
\end{table}

Stage 1 is measured in minutes: the alert fires when the survey's astrometric
solution is good enough for a follow-up to be scheduled. Stage 2 is measured
in weeks for a short-arc object and in years for one that will return, and it
is the stage at which the astrometric term dominates for most of the
relevant horizon. Stage 3 has no fixed clock at all, because it has no
trigger: it happens when an observer with an idle telescope and an interest
happens to be pointed at the right object. Stage 4 and stage 5 are outside
this book's scope, which is deliberate; the endpoint here is accountability,
not deflection.

\section{The trigger that does not exist}

The most consequential sentence in the 2025 NASA documentation on the
2024~YR4 episode is a negative one. The planetary-defence coordination
office and the international warning network were recorded as discussing the
creation of an automated system, and the documented position was that
\textbf{no automated system currently alerts physical characterization
sites}.

The operational consequence was observed in real time. Observers at the
Infrared Telescope Facility had not learned of the object's Torino Level 1
status when they were observing it. That is, in the most literal sense, a
failure to detect: the object was detected, imaged and measured, and the
people doing so were not in the loop that would have told them the
measurement mattered.

The same documentation records that an automated alert was exercised
successfully for a Torino Level 1 event --- 2025~FA22, on 25 April 2025 ---
which establishes that the trigger is operationally simple rather than
technically difficult. The 2024~YR4 episode was therefore not a sensing
failure and not a data-processing failure. It was a routing failure.

\begin{seamnote}{A sensor does not fix a routing failure}
The owner's original concept carried a beam to Earth, and the instinct behind
it was good: a link from the risk list to an instrument that can respond. The
2024~YR4 record says the link was the missing piece, and that an automated
alert to a characterisation site is the concrete form that the link should
take. This book therefore treats the beam as an alerting\idx{Alerting} and scheduling path
rather than as a ranging channel, because the ranging channel turns out to be
the weaker claim of the two.
\end{seamnote}

\section{Where the Sun-ward blind zone bites}

Stage 1 is not equally accessible from the ground. A ground telescope cannot
observe close to the Sun: the atmospheric scattering and the sky background
rise, and the solar avoidance constraints tighten. Objects arriving from the
Sun-ward direction are therefore found late or not at all, and this is not a
hypothetical concern. The 2026 study of near-Earth-object telescopes for
extreme environments names the historical failures directly: 2008~TC3 and
2019~OK, and ``multiple other NEAs detected only days before their closest
approach to Earth''.

The geometry of that disadvantage can be inverted from the published table in
that study. For a relative velocity of 20~kilometres per second, a 140-metre
object becomes detectable from a lunar observatory at 0.10 astronomical units
and about nine days earlier than from the ground; a 30-metre object at
0.05~astronomical units and four to five days earlier; and a one-kilometre
object at 0.50~astronomical units and about six weeks earlier. Working
backwards from the lead times gives the ground limits the table implies:
about 0.20 astronomical units for the 140-metre class, about 0.10 for the
30-metre class and about 0.99 for the kilometre class.

Those are the numbers in \figref{fig:sunward} of Chapter 26, and they are the
discovery-side case for a lunar platform. The lead times are modest, and the
book does not inflate them. Their significance is not their length but their
direction: they are the only entry in the table that applies to objects the
ground structurally cannot reach.

\begin{keytakeaways}
\item The two mandates are separately funded and separately staffed, and only
one of them has a trigger.
\item The documented 2024 YR4 failure was a routing failure: the observers
were not told the object had reached Level 1.
\item An automated alert to a characterisation site was exercised
successfully for 2025 FA22 on 25 April 2025, so the missing link is
operationally simple.
\item The stage that actually binds is characterisation, and the Sun-ward
blind zone is the one part of discovery that a lunar vantage reaches and the
ground does not.
\end{keytakeaways}
